\section{Marco Te\'orico}
\label{sec:marco_teorico}

\subsection{Riesgo como producto de peligro, exposici\'on y vulnerabilidad}

La gesti\'on de riesgo de desastres descompone convencionalmente el riesgo en tres factores multiplicativos: \emph{peligro} (la probabilidad de que ocurra un evento f\'isico detonante, p. ej. un sismo o una tormenta extrema), \emph{exposici\'on} (qu\'e o qui\'en est\'a en la trayectoria de ese evento) y \emph{vulnerabilidad} (qu\'e tan susceptible es el elemento expuesto a sufrir da\~no dado el evento). Aplicado a una presa, el peligro es externo al sitio (sismicidad regional, r\'egimen de lluvias), la exposici\'on es la poblaci\'on y la infraestructura aguas abajo, y la vulnerabilidad es intr\'inseca a la estructura misma --su geometr\'ia, edad, material y estado de conservaci\'on--. Los tres factores son necesarios: una presa en zona de bajo peligro s\'ismico pero estructuralmente deteriorada puede representar m\'as riesgo que una presa robusta en zona de alto peligro.

\subsection{Metodolog\'ias de \emph{screening} para portafolios grandes de presas}

Cuando un operador es responsable de cientos o miles de presas, evaluar cada una con un an\'alisis estructural detallado es inviable en tiempo y presupuesto; de ah\'i que la literatura en ingenier\'ia s\'ismica haya desarrollado metodolog\'ias de evaluaci\'on r\'apida (\emph{screening}) para priorizar qu\'e presas ameritan un an\'alisis m\'as profundo primero. Hariri-Ardebili y Nuss \cite{hariri2018prioritization} proponen un \'indice sem\'i-cuantitativo de priorizaci\'on s\'ismica que combina par\'ametros de la presa, del embalse, y de riesgo aguas abajo. Un enfoque m\'as reciente, orientado a fases de preparaci\'on y respuesta a emergencias, propone una herramienta multiescala que arranca con an\'alisis simplificados de cuerpo r\'igido en dos dimensiones y solo escala a modelos no lineales cuando el screening inicial lo justifica, precisamente para mantener el proceso \'agil y econ\'omico sobre portafolios grandes \cite{natural2023multiscale}. Ambos enfoques, sin embargo, presuponen que el operador ya cuenta con los atributos b\'asicos de cada presa del portafolio --altura, geometr\'ia, tipo de cortina-- capturados en un inventario propio; el problema que resuelven es c\'omo \emph{ordenar} un an\'alisis m\'as profundo, no c\'omo operar cuando esos atributos b\'asicos no est\'an disponibles en absoluto para la mayor parte del portafolio.

\subsection{El vac\'io de datos como el problema previo a resolver en M\'exico}

Ese es precisamente el punto en el que la situaci\'on de M\'exico se aparta de los casos de estudio t\'ipicos en la literatura. No se trata de decidir qu\'e nivel de an\'alisis estructural aplicar --pseudo-est\'atico, din\'amico lineal, no lineal-- sino de que los datos de entrada m\'as b\'asicos para \emph{cualquiera} de esos niveles (altura de cortina, a\~no de construcci\'on, capacidad del vertedor, estado f\'isico y fecha de \'ultima inspecci\'on) no est\'an publicados para la mayor parte de las presas del pa\'is, y en particular no lo est\'an para el subconjunto de \(\sim\)1{,}000 presas que CONAGUA reconoce no haber inspeccionado recientemente \cite{infobae2020presas} --el subconjunto donde un \'indice de priorizaci\'on ser\'ia m\'as \'util--. Esta situaci\'on no es exclusiva de M\'exico: la literatura reconoce que los vac\'ios de datos espaciales para presas peque\~nas y medianas persisten en regiones en desarrollo, con omisiones y actualizaciones retrasadas que dificultan estimar el riesgo de desastre asociado a esta infraestructura \cite{natural2023multiscale}.

\subsection{Contribuci\'on de este trabajo}

Frente a ese vac\'io, este trabajo no intenta emular un an\'alisis estructural que los datos disponibles no permiten. En su lugar, construye la mitad del problema de riesgo que s\'i es calculable enteramente con datos abiertos de terceros --peligro s\'ismico regional, r\'egimen de lluvia extrema, y exposici\'on poblacional-- como un \'indice de \emph{priorizaci\'on de inspecci\'on}, dejando expl\'icito qu\'e falta (la vulnerabilidad estructural) y d\'onde se est\'a gestionando (una solicitud de transparencia gubernamental, detallada en la Secci\'on \ref{sec:metodologia}). Esto sigue el esp\'iritu del \emph{screening} en dos etapas de \cite{natural2023multiscale} --evaluaci\'on r\'apida primero, an\'alisis detallado despu\'es solo donde se justifique-- pero aplicado a la etapa previa que la literatura revisada no aborda expl\'icitamente: qu\'e hacer cuando ni siquiera el inventario b\'asico es p\'ublico.
